\subsection{Simmetria di gauge}
Abbiamo visto diversi tipi di simmetrie: Simmetrie per traslazione (anche temporale), rotazione. Simmetrie discrete come C, P, T anche perchè la violazione di queste simmetrie non è per niente ovvia ma anche delle applicazioni per poter giustificare che il tempo abbia una freccia e della simmetria tra materia e antimateria. Simmetria nello spazio interno delle varaibili (isospin, SU(3)). Ora si parla di simmetrie di Gauge. Compaioni in maniera naturale nell'elettromagnetismo classico però acquistano un significato più profondo in meccanica quantistica relativistica: Estensione delle simmetrie nello spazio interno delle particelle e introducono in maniera univoca interazioni collegate con queste simmetrie interne. Introduzione del modello standard: ci sonotre gruppi di simmetria chiamati: $U(1)_Y$ di ipercarica, $SU(2)_L$ indicato left perchè nelle interazioni deboli abbiamo una certa asimmetria di elicità (le particelle hanno elicità negativa mentre le antiparticelle positiva) si parla anche di chiralità che è la sua estensione. Mentre $SU(3)_C$ di colore. Le simmetrie $U(1)_Y$ e $SU(2)_L$ sono in realtà simmetrie che non sono esatte a basse energie a causa del processo di rottura spontanea della simmetria che mi separa fisicamente le interazione elettromagnetiche con quelle deboli anche se posso creare una teoria combinata di queste (interazioni elettrodeboli). 

\subsection{Il Modello Standard}
%IMMAGINE ewe 
\begin{figure}[htp!]
  \centering
  \includegraphics[scale=0.3]{Lezione19/dsa.png} 
\end{figure}
Il Modello Standard è basato su tre simmetrie di gauge: per le interazioni forti: gruppo SU(3) di colore, grado di libertà interno dei quark. Quindi ho 8 generatori e quindi 8 gluoni. Particelle senza massa. Questi gluoni sono confinati all'interno dei barioni ($qqq$) e dei mesoni ($q\Bar{q}$). Le interazioni elettrodeboli sono prodotti da due gruppi: U(1) di ipercarica ha un generatore mentre $SU(2)$ di isospin debole corrisponde a 3 generatori $Z$ $W^+$ $W^-$. Essendo un gruppo SU(2) organizzate in doppietti tutti i fermioni del modello standard sono organizzati in doppietti. I quark: $Y=1/3$. I leptoni invece $Y=-1$ e vale sempre la legge di Gell-Mann Nishijima:
\begin{equation*}
    Q=\frac{Y}{2}+T_3
\end{equation*}
La cosa importante è che mentre il fotone ha massa nulla, la massa di $W$ e $Z$ è molto grande infatti $W$ ha circa $80$ volte la massa del protone. Particelle che hanno una massa di un nucleo di medie dimensioni. Quello che succede però è che le simmetrie di gauge vogliono particelle di massa nulla. Questa cosa sembra essere contro il voler spiegare le interazioni tramite le simmetrie di gauge. Per fare questo quello che abbiamo è che le interazioni deboli amano violare tutto. Violano C violato P di conseguenza violano anche T e anche la simmetria di Gauge con il meccanismo di rottora spontanea della simmetria con un'altra particella che nel Modello Standard dobbiamo attaccare come chiave di volta per far tornare tutti i conti che è il \textbf{Bosone di Higgs}.
\subsection{Proprietà di W e Z}
%IMMAGINE ewe 
\begin{figure}[htp!]
  \centering
  \includegraphics[scale=0.2]{Lezione19/yu.jpg} 
\end{figure}
Sono particelle di massa grande quindi possono decadere e lo fanno molto velocemente. Se vado a vedere la larghezza di $W$ effettivamente ho una larghezza maggiore della tipica larghezza delle interazioni forti. Infatti avevamo visto che "durano" $10^{-24}\,s$ o $10^{-25}\,s$ qui invece se traduciamo la larghezza in un tempo siamo nell'ordine di $10^{-26}\,s$ quindi durano molto meno. Non riescono neanche ad attraversare un nucleo prima di decadere. Tutte le particelle mediatrici hanno spin $=1$ con carica $\pm1$. La $Z$ si può vedere come la componente di tripletto di isospin con $T_3=0$. Questi sono i due primi eventi osservati.  Il $W$ veniva prodotto e decadeva in un muone e in un neutrino mentre la $Z$ in due elettroni. La cosa importante che la $Z$ è stata una delle particelle più difficili da rilevare e da sviluppare teoricamente. Questo perchè per $W$ abbiamo tutti i decadimenti adronici che ci permettono di misurare la costante di fermi in diversi processi. Siccome la $Z$ non cambia il sapore delle particelle non può mediare decadimenti e non potendo mediare decadimenti di fatto non si riusciva a studiare fenomenologia se non con processi indiretti.

Prima di passare alla discussione del Bosone di Higgs bisogna fissare per bene la finestrella del modello standard perchè è davvero ricco di informazioni! I quarks sono le particelle che interagiscono forte, i leptoni invece interagiscono solo attraverso interazioni deboli ed elettromagnetiche (tranne i neutrini che interagiscono solo deboli perchè non hanno la carica). All'interno delle Force particles c'è anche l'informazione del range delle interazioni deboli che spiega la teoria di fermi del decadimento beta. 
\subsection{Rottura spontanea di simmetria}
Abbiamo visto come la simmetria di gauge diventa un meccanismo che accomuna i campi collegati alle interazioni con i campi collegati con le particelle interagenti e quindi con le trasformazioni e i gradi di libertà interni di queste particelle interagenti. Questa simmetria di gauge si usano per analizzare tutte e tre le interazioni ma questo prevede che tutte le particelle di interazioni abbiano massa nulla. Il fotone ha massa nulla il gluone ha massa nulla quindi le interazioni forti ed elettromagnetiche sono a posto e soddisfano questa richiesta esplicitamente. Mentre i bosoni vettori delle interazioni deboli sono molto massive. La questione è come posso conciliare le due cose? Da una parte ho la simmetria di gauge che mi piace molto perchè collega le trasformazioni di simmetria a interazioni in maniera univoca con il fatto che ho bosoni vettori con una massa. Il modo di farlo è introdurre la \textbf{rottura spontanea della simmetria} significa che la mia teoria è formalmente simmetrica ma lo stato fondamentale no. Sembra una contraddizione. Questo viene introdotto attraverso il meccanismo di Higgs che aggiunge una particella in più che serve a chiudere questo modello consistente che descrive in maniera molto dettagliata le interazioni. 

La rottura di simmetria è un fenomeno in cui le equazioni che descrivono un sistema mostrano una certa simmetria ma lo stato fondamentale viola questa simmetria. Di questi esempi ci sono anche in meccanica classica ad esempio i materiali ferromagnetici. In questo caso i momenti magnetici dei vari atomi tendono ad allinearsi perchè c'è un'interazione tra i momenti magnetici favorendo energeticmamente il caso in cui ho i momenti magnetici allineati. L'hamiltoniana che descrive questo processo è esplicitamente invariante per rotazione:
\begin{equation*}
    H=-k\sum_{i\ne j}\Vec{\mu}_i\cdot\Vec{\mu}_j
\end{equation*}
Però ad un certo instante i momenti magnetici si allineeranno in una direzione comune. Questo (nello stato fondamentale e quindi di minima energia) produrrà una direzione "privilegiata" non avendo più invarianza per rotazione. La risposta dipende anche dalla temperatura.

Quindi a parte la dimostrazione io ho che la simmetria di gauge è sempre conservata perchè l'equazione originale da dove siamo partiti è rispettata ma se sviluppiamo l'equazione in un punto di minimo del potenziale vedo che l'apparenza fisica di questo è che la simmetria è rotta e compare un termine di accoppiamento campo-potenziale e questo termine mi genera la massa dei bosoni vettori. La rottura spontanea di simmetria fa acquisire un termine di massa ai campi mediatori delle forze. Questo è il modello con cui diamo una massa ai bosoni $W$ e $Z$. Chiaramente per fare questo è necessario un campo e quindi una particella che abbia questa proprietà di rottura della simmetria nel suo stato fondamentale.
\subsection{Bosone di Higgs}
Quando descritto in maniera estremamente semplificata è il concetto alla base del meccanismo di Higgs. Questa particella quando è stato fatto il modello stranamente si conoscevano altre proprietà ma non la massa. Ovviamente doveva avere spin $=0$ perchè le equazioni di Klein-Gordon descrivono una particella senza spin. Ma soprattuto non so quanto vale la sua massa ma conosco il suo valore di aspettazione $v=246\,GeV$. Nel 2012 questa particella è stata finalmente osservata. L'ipotesi di Higgs è del 1964, il modello standard del 1972-1973 e ci sono voluti quindi 40 anni prima di arrivare alla conferma piena. $m_H=125\,GeV$, $Q=0$, $S=0$. Quindi si mette in una categoria a se stessa. Questa è l'unica particella elementare con spin nullo. Questo bosone di Higgs è quello che fornisce la massa a $W$ e $Z$.